\section*{Chapter \ref{ch:hf:queue-reactive-models-and-large-tick-assets} --- Queue-Reactive Models and Large-Tick Assets}

\begin{solution}{exo:hf:queue-reactive-models-and-large-tick-assets:1}
$290/200=1.45$ lots a second.
\end{solution}

\begin{solution}{exo:hf:queue-reactive-models-and-large-tick-assets:2}
$s^b=3/15=0.2$: below 0.3, so it does not join the bid; the ask's share is 0.8, so it joins the ask.
\end{solution}

\begin{solution}{exo:hf:queue-reactive-models-and-large-tick-assets:3}
$0.392\times3\,400\times\$0.01=\$13.33$.
\end{solution}

\begin{solution}{exo:hf:queue-reactive-models-and-large-tick-assets:4}
A short opposite queue empties soon, and the price then moves away from the resting order before it is reached: the order is left unfilled.
With a long opposite queue the price is unlikely to move away, so even the back is eventually filled.
\end{solution}

\begin{solution}{exo:hf:queue-reactive-models-and-large-tick-assets:5}
Joining on a favourable share selects good states only at the moment of joining. Orders that stay are then filled mostly after the share has
turned against them, when their queue is being eaten ahead of a price move: the rule keeps precisely the orders whose fills are adverse
($-0.117$ ticks a share against 0.026 for always joining).
\end{solution}

\begin{solution}{exo:hf:queue-reactive-models-and-large-tick-assets:6}
A higher threshold keeps only the best states, so each fill is better, but it quotes less often and fills fewer shares; the edge per session is
the product, which peaks at 0.3 (\$13.33) while the mark-out per share peaks at 0.4 (0.484).
\end{solution}

\begin{solution}{exo:hf:queue-reactive-models-and-large-tick-assets:7}
Fill probability 0.825, value 0.594 ticks a share, adverse move $-0.22$: after its fills the mid moves in its favour on average, because the
short ask queue tends to empty and the price to rise.
\end{solution}

\begin{solution}{exo:hf:queue-reactive-models-and-large-tick-assets:8}
The model has no informed flow: its only adverse selection is mechanical. The simulated market's informed traders make fills worse exactly when
queues are being eaten, which the model's values do not include; chapter 1's quoter, always in the queues, lost its spread. Measure mark-outs of
one's own fills and quote only where they are positive.
\end{solution}

\begin{solution}{pb:hf:queue-reactive-models-and-large-tick-assets:1}
\begin{enumerate}
\item A spread almost always one tick with long best queues; the market maker controls which queues to be in and when to leave them, not the price.
\item One tick 96.3\% of the time; median best queue 17 lots.
\item Market orders here are too small to exhaust a best queue, so the next price is almost never reached.
\item See \cref{def:hf:queue-reactive-models-and-large-tick-assets:qr}.
\item \Cref{met:hf:queue-reactive-models-and-large-tick-assets:estimate}; limit orders between 1.3 and 2.3 lots a second with no trend, cancellations rising
from 0.34 to about 1.5, market orders falling from 0.55 to about 0.3.
\item Each resting order cancels at its own rate, so a longer queue loses lots faster.
\item One lot in the bid queue with $n$ ahead; joins, cancellations and market orders at the fitted rates; filled when reached, abandoned if the ask
empties or after 60 seconds; after a fill, marked at the mid ten seconds later: $0.5$ tick plus the mid's move.
\item 0.60 (front) and 0.32 (back) at 10 against 2; 0.12 for both at 2 against 20.
\item The queue is about to be eaten: every place is filled, followed by the price move.
\item The model has no informed flow; the information in the flow is omitted.
\item See \cref{def:hf:queue-reactive-models-and-large-tick-assets:rules}.
\item Always joining 0.026, joining and staying $-0.117$, joining and leaving 0.392 ticks a share; in the model, the front of a ten-lot queue against
a two-lot ask is worth 0.60 ticks and the back 0.32.
\item It marks each session's end inventory, whose moves are large against the rules' differences over six sessions.
\item When the other side is about to take the queue, the front is filled first and the move follows.
\item 0.3 for the edge per session; the mark-out keeps rising to 0.4 because stricter rules keep better but fewer fills.
\item Always 365, stay 245, leave 468, leave unless at the front 448: leaving means cancelling and rejoining as the book moves.
\item Fit the fills' mark-out, by state, to the book's features at the fill and add it to the payoff, or add informed market orders whose rate
depends on a hidden efficient price.
\item Under pro-rata allocation the place in the queue does not matter, size does (chapter 15); leaving costs no place.
\item The spread is several ticks, so the choice of price reappears (chapters 3--4) and queues are short; the rules matter less.
\item Leaving well.
\end{enumerate}
\end{solution}

\begin{solution}{iq:hf:queue-reactive-models-and-large-tick-assets:1}
The ask is about to empty and the price to move up: a bid at the front will probably be filled well, and the move follows in its favour. Stay; the
danger is the reverse, a short bid queue.
\iqlookfor{reading the other side's queue, and that the front is not always safe.}
\end{solution}

\begin{solution}{iq:hf:queue-reactive-models-and-large-tick-assets:2}
Record each order's place at every event and its outcome (fill, mark-out at a horizon, or cancellation); estimate the expected mark-out and fill
probability by place and by the queues' sizes, out of sample; the difference between places is the queue value.
\iqlookfor{outcomes conditioned on place, and selection bias of which orders stay.}
\end{solution}

\begin{solution}{iq:hf:queue-reactive-models-and-large-tick-assets:3}
An order far back is reached only when most of the queue has been eaten, which is when the price is about to move through it: higher fill
probability late in a queue's life comes with worse fills.
\iqlookfor{the conditional nature of late fills.}
\end{solution}

\begin{solution}{iq:hf:queue-reactive-models-and-large-tick-assets:4}
From an order-by-order feed: at its acknowledgement the order joins behind every order resting at its price; each execution or cancellation of an
order ahead moves it forward. Hidden orders, lost or reordered packets and venue-specific priority rules (size changes, replaces) break the count.
\iqlookfor{order-by-order tracking and its failure modes.}
\end{solution}

\begin{solution}{iq:hf:queue-reactive-models-and-large-tick-assets:5}
Every cancel and re-entry sends the order to the back of the queue: in a one-tick stock the place is most of the order's value, so churning gives it
away, besides adding messages.
\iqlookfor{priority lost on replace.}
\end{solution}

\begin{solution}{iq:hf:queue-reactive-models-and-large-tick-assets:6}
With $k$ ahead the events among them occur at rate $\mu+k\nu$; the next is a market order with probability $\mu/(\mu+n\nu)$. Each event removes one
order ahead, so the expected time until none is left is $\sum_{k=1}^n1/(\mu+k\nu)$.
\iqlookfor{competing exponential clocks and summing the stages.}
\end{solution}
